%% kapa: exam=1 round=29 year=2018 period=2 form=A subject=appraisal_law subject_ko=감정평가관계법규 seq=23 number=23
%% source: Q-Net 감정평가사 기출문제 (cst00309) artlSeq=5202808 file=extracted/1008/5202808/2018년 제29회 감정평가사 1차시험 기출문제/2018년 제29회 감정평가사 1차시험 기출문제/2교시 A형.pdf
%% conversion: kordoc 4.24.51 -> tools/gen1.py (auto); review=proofread-sonnet; flags=none
\begin{question}{23}
국유재산법상 국유재산 관리ㆍ처분의 기본원칙으로 명시되어 있는 것은?
\begin{choices}
  \item 수익과 손실이 균형을 이룰 것
  \item 보존가치와 활용가치를 고려할 것
  \item 투명하고 효율적인 절차를 따를 것
  \item 지속가능한 미래의 가치와 비용을 고려할 것
  \item 국유재산이 소재한 지방자치단체의 이익에 부합되도록 할 것
\end{choices}
\end{question}
\answer{3}
